\chapter{Introduction}
\label{chap:introduction}

Public-key cryptography underpins virtually all secure communication on the internet. VPN tunnels, TLS handshakes, SSH connections: they all depend on the computational hardness of problems like integer factorization and the elliptic-curve discrete logarithm, problems that classical computers cannot efficiently solve at cryptographic parameter sizes. This assumption has held for decades, but it will not hold forever.

Peter Shor's 1994 polynomial-time quantum algorithm~\cite{shor1994algorithms} solves both integer factorization and the discrete logarithm problem, which means a sufficiently large fault-tolerant quantum computer would break RSA and elliptic-curve cryptography in polynomial time. The security community has been aware of this since the 1990s, but for most of that time quantum hardware was too primitive to pose a real threat. That is changing. Current estimates put the cost of factoring RSA-2048 at roughly millions of physical qubits with today's error-correction techniques~\cite{gidney2021factor}, and hardware roadmaps from IBM, Google, and others suggest that scale could be reached within the next decade or two.

In response, NIST launched its Post-Quantum Cryptography (PQC) Standardization Project in 2016, finalizing the first round of standards in August 2024: FIPS 203 (ML-KEM, for key encapsulation)~\cite{fips203}, FIPS 204 (ML-DSA, for digital signatures), and FIPS 205 (SLH-DSA, hash-based signatures). For VPN key exchange, ML-KEM is the directly relevant standard. NIST has also set a hard timeline: classical key exchange algorithms are to be deprecated by 2030 and disallowed entirely by 2035~\cite{nistir8547}.

WireGuard is among the most important migration targets. Since Donenfeld introduced it in 2017~\cite{donenfeld2017wireguard}, it has largely replaced IPsec and OpenVPN as the go-to VPN protocol, thanks to its small codebase, strong performance, and easy auditability. It ships in the Linux kernel, runs on every major operating system, and is deployed at scale by Cloudflare, Mullvad, NordVPN, Tailscale, and many enterprises. But WireGuard's cryptographic design is deliberately opinionated: a fixed cipher suite with no negotiation. This rigidity is a security advantage (cipher agility has been a recurring source of vulnerabilities in TLS and IPsec), but it makes migration harder, since there is no built-in way to swap in a different key exchange algorithm.

WireGuard does, however, provide two natural handles for post-quantum augmentation. The first is the preshared-key (PSK) slot, a 256-bit symmetric key that is mixed into the handshake's key derivation chain. If a post-quantum KEM is used out-of-band to establish this PSK, the resulting handshake is quantum-resistant against passive adversaries, at the cost of requiring separate key distribution infrastructure. The second approach is to modify the Noise IK handshake itself, embedding ML-KEM operations alongside the existing X25519 exchange so that both shared secrets contribute to the final session key in every handshake. The latter is what we call the full hybrid approach.

This thesis makes four contributions. It implements a PSK-slot baseline: a small CLI tool that performs ML-KEM-768 key encapsulation between two peers out-of-band and feeds the resulting 256-bit key into BoringTun's~\cite{boringtun} preshared\_key field. It then implements the full hybrid: a modified BoringTun whose handshake module runs both X25519 and ML-KEM-768 in every handshake, mixing both shared secrets into the Noise chaining key via HMAC-BLAKE2s, with no changes to the data-transport phase. Both implementations are evaluated on Apple Silicon (M4) and Raspberry Pi 5 hardware, measuring handshake overhead and packet sizes. Finally, a security analysis compares the two approaches and discusses where each one fits in practice.

The rest of the thesis is structured as follows. Chapter~\ref{chap:background} covers the relevant background: the quantum threat, WireGuard, the Noise framework, and ML-KEM. Chapter~\ref{chap:related_work} surveys prior work on post-quantum WireGuard. Chapter~\ref{chap:problem} states the research question. Chapters~\ref{chap:design} and~\ref{chap:implementation} describe the design and implementation, respectively. Chapter~\ref{chap:security} analyzes security properties. Chapter~\ref{chap:evaluation} reports the evaluation results. Chapter~\ref{chap:discussion} discusses tradeoffs, deployment considerations, and future work, and Chapter~\ref{chap:conclusions} concludes.

\section*{Declaration of Generative AI and AI-Assisted Technologies in the Writing Process}

During the preparation of this work I used Claude (Anthropic's large language model) via the Claude Code CLI tool for these purposes: as a coding assistant while implementing the hybrid handshake and PSK-slot tool in Rust, for designing and interpreting the benchmarks, and as an editor to tighten the thesis prose and catch places where the written descriptions had drifted from what the code actually does. After using this tool, I reviewed and edited the content as needed and take full responsibility for the content of the thesis.
